\subsection{A maximal sporadic group}
Put $F=\mathbb Q(z)$, where $z^2+z+1=0$, and put $X=X_1(3,15)$.
We determine the decic curves whose torsion contains $C_3\times C_{15}$.
The following saturation calculation also strengthens the
Jacob-ian input used for the octic case.

\begin{lemma}\label{lem:315-cuspidal-generation}
The group $J(F)$, where $J=\operatorname{Jac}(X)$, is generated by
$F$-rational cuspidal divisors and has invariant factors
$[3,3,24,96,96]$. Reduction at the prime above $2$ is injective on $J(F)$.
\end{lemma}
\begin{proof}
The Jacobian has rank zero over $F$ by \cite[Theorem~4.1]{DS}.
Let $C\subseteq J(F)$ be generated by $F$-rational cuspidal divisors.
The accepted cusp calculation, repeated at $4$ and $7$, gives
\[
 C\simeq[3,3,24,96,96],\qquad \#C=1990656.
\]
The prime $7$ splits in $F$, and reduction there is injective on torsion.
The Jacobian orders at the split primes $7,13,19$ are
$149299200$, $12230590464$, and $247669456896$.
Their gcd is $3\#C$. Consequently $[J(F):C]\mid3$, and their
$2$-primary subgroups agree.

At the prime above $2$ the complete divisor class group has degree-zero
invariant factors $[3,3,3,24,96,96]$. Its $3$-primary subgroup
is therefore $\mathbb F_3^6$.
The classes of all closed cusps over $\mathbb F_4$, after projecting
to the $3$-primary part, span this whole space. This rank-$6$
assertion is checked by the exported cusp class vectors.

All cusps split in $L=F(\zeta_5)$, with
$\operatorname{Gal}(L/F)\simeq C_4$. Frobenius at $4$ has order $2$.
Thus every closed cuspidal divisor over $\mathbb F_4$ lifts to a
cuspidal divisor defined over the quadratic extension
$L_0/F$ fixed by this Frobenius. Manin--Drinfeld makes its
class, relative to a rational cusp, torsion.
Reduction is injective on its $3$-primary part. Hence every
point of $J(F)[3^\infty]$ has order dividing $3$ and is a
linear combination of these lifted $3$-primary cuspidal classes
in $J(L_0)$.

Let $b$ be such an $F$-rational point. Averaging that expression
over $\operatorname{Gal}(L_0/F)$ gives $2b$ as a combination of
$F$-rational cuspidal classes. Since $b$ has order dividing $3$,
it follows that $b\in C$. This removes the possible index $3$,
and proves $J(F)=C$.
The cuspidal image at $4$ has order $\#C$, so reduction is injective.
As an additional check, the involution on the six-dimensional local
$3$-space has fixed space of dimension $5$, exactly the image of $C[3]$.
\end{proof}

\begin{lemma}\label{lem:315-degree5-classes}
The curve $X$ has $F$-gonality greater than $5$.
There are at most $48$ classes of effective $F$-rational degree-$5$
divisors whose support contains no cusp, after removing the known
reducible classes.
\end{lemma}
\begin{proof}
We identify global cusp divisors at $4$ and $7$ using the valuations of
\[
 a_3,\ a_1^3-27a_3,\ x(kQ+jP),\ y(kQ+jP),\ x(kQ)-x(kQ+P),
 \qquad1\leq k\leq7,\quad0\leq j\leq2,
\]
on the normalized universal curve
$y^2+a_1xy+a_3y=x^3$, with $P=(0,0)$ and $Q$ of order $15$.
These are integral units on the noncuspidal locus away from $15$.
Their valuation vectors distinguish the twenty global cusp blocks,
of degrees $1^{16},4^4$. At $4$, each quartic block splits into two
quadratic closed cusps. The common free cusp maps have equal kernels.

The independent Frobenius census gives the following numbers of
noncuspidal closed places:
\[
\begin{array}{c|rrrrr}
q\backslash d&1&2&3&4&5\\\hline
4&0&0&0&102&192\\
7&0&24&144&396&3312
\end{array}
\]
At $4$ we enumerate all $24432$ effective degree-$5$ divisors:
$22608$ purely cuspidal divisors, $102\cdot16$ quartic-plus-cusp
patterns, and $192$ quintic places.
At $7$ the noncuspidal patterns have counts
\[
 19584,\ 19584,\ 6336,\ 3312,\ 4800,\ 3456
\]
for a single place of degree $2,3,4,5$ with its full cusp complement,
and for the patterns $2+2+1$ and $2+3$, respectively.
Their total is $57072$. Together with the $15568$ purely cuspidal
divisors, all $72640$ effective degree-$5$ divisors have distinct
classes over $\mathbb F_7$.
Thus there is no function of degree at most $5$ over $\mathbb F_7$.
Specialization gives the same lower bound over $F$.
Forgetting the independent $3$-torsion point gives a degree-$3$ map
to $X_1(15)$, which has genus one and a rational cusp.
Composing with a degree-$2$ function shows that the $F$-gonality is exactly $6$.

A noncuspidal global divisor cannot have purely cuspidal reduction
at $7$. Indeed every such local cusp divisor lifts over $F$, since
Frobenius at $7$ generates $\operatorname{Gal}(L/F)$.
Odd-prime injectivity on torsion and the gonality bound would make
the global divisor equal to this cuspidal lift.

By Lemma~\ref{lem:315-cuspidal-generation}, the two reductions of a
global degree-$5$ class must lie in the exact diagonal image of $C$.
The integer Hermite-normal-form replay gives precisely $240$ possible
classes. Of these, $192$ have known global representatives obtained
by adding each rational cusp to the twelve degree-$4$ divisors from
the octic calculation. Each representative contains a cusp.
This leaves at most $48$ classes with no cusp in the support.
\end{proof}

\begin{proposition}\label{prop:decic315}
The group $C_3\times C_{15}$ occurs as full torsion over a decic field.
Every elliptic curve over a decic field whose torsion contains this
group has full torsion $C_3\times C_{15}$.
There are exactly five resulting $\overline{\mathbb Q}$-isomorphism
classes of elliptic curves.
\end{proposition}
\begin{proof}
Let $K=F(w)$, where
\[
 w^5+13(z+1)w^4-(24z+22)w^3+4(z+1)w^2+(z+4)w+z+1=0.
\]
The polynomial is irreducible over $F$. Put
\[
 q=\frac{(293z+23)w^2+(105-238z)w+357-28z}
         {217w^2+(603z+271)w+283+86z}.
\]
On the curve
\[
 E:y^2+((1+q)w+2-q)xy+(qw^2+(1-q)w)y=x^3,
\]
the points $P=(0,0)$ and $Q=(-w,w^2)$ have exact orders $3$ and $15$,
with $e_3(P,5Q)=z^2$. At the unramified prime above $43$ given by
$z\mapsto6$, $w\mapsto28$, the curve reduces to $y^2+36xy+24y=x^3$, with group
$[3,15]$. Consequently $E(K)_{\rm tors}=C_3\times C_{15}$.
All identities and the reduction count are checked independently
in Magma and Sage.

The $j$-invariant has degree $5$ over $\mathbb Q$ and is not integral,
so $E$ is non-CM. In particular $K=F(j(E))$.
The $F$-defined group
\[
 G=(\operatorname{SL}_2(\mathbb F_3)\times\mathbb F_5^\times)
       /\langle(-I,-1)\rangle
\]
has order $48$ and changes only the level structure.
Its orbit on the resulting closed degree-$5$ point has size $48$:
an automorphism preserving that closed point induces a field
 automorphism fixing $j(E)$, hence is the identity on $K$;
the non-CM fine moduli interpretation then makes it trivial in $G$.
The gonality bound makes the $48$ divisor classes distinct.
They therefore exhaust the $48$ remaining classes of
Lemma~\ref{lem:315-degree5-classes}.

Now take any decic curve containing $C_3\times C_{15}$.
Its field contains $F$, and its moduli trace is a degree-$5$
divisor over $F$ supported away from the cusps. It cannot belong
to one of the $192$ reducible classes, by the gonality bound.
It belongs to one of the $48$ classes just constructed, and
uniqueness of effective divisors in that class identifies its
fine moduli point with the displayed example after changing
level structure and conjugating the field. Its full torsion is
therefore again $C_3\times C_{15}$.
The five conjugates of $j(E)$ give exactly five geometric
isomorphism classes.
\end{proof}

\begin{corollary}\label{cor:decic315-maximal}
Neither $C_3\times C_{30}$ nor $C_3\times C_{45}$ occurs over a decic
field, and no strict supergroup of $C_3\times C_{15}$ occurs.
\end{corollary}

The exact point can be verified from its displayed equations alone.
The maximality proof uses the serialized finite class groups,
Frobenius polynomials, cusp vectors, and complete local divisor lists.
The separate Sage replay checks the census, Galois saturation,
all $72640$ distinct classes at $7$, and all diagonal matches by integer
Hermite normal form. The initial three-multiple sieve and the failed
dyadic diagnostic are discovery records and are not substitutes for
these exact checks.
